\chapter{术语表：一份受控词汇的五条概念路径}
\setcounter{section}{-1}

\begin{noteblock}
Semantica 绑定： 本附录是 185 个词条及五条概念路径的人可读解释；机器投影的
唯一位置是 \nolinkurl{semantica.chapter_packages.vol2.normative}，其中 Part 1 资产登记为
\nolinkurl{part1-vocabulary} 与 \nolinkurl{part1-cards}。该 package 当前为 \nolinkurl{partial}、release
\nolinkurl{blocked}；它保存的是本书发布的工程释义、来源坐标和教学映射，不是 ISO 原文、
授权译文、官方解释或合规证据。本目录没有 RDF 锚点或查询副本。
\end{noteblock}

\begin{noteblock}
本附录以 ISO 26262-1:2018 第 3 章（§3.1–§3.185，1-based 物理页 9–36）为唯一受控词条集，
按本附录建立的五个概念群组织正文（与第 2 章的术语铺垫相衔接），字母序只作为末尾的入口索引。
附录不建立第二套术语体系：每个条目的“一句白话”是教学重述，不是授权译文，
更不能替代原文词条在合同、审核或评估场合的效力。

\textbf{与全书二十章的衔接}：本书前十章现为叙事重写，方法表语义、独立性等级记号、FMEDA 与失效分类等器具性内容已集中由附录承载；正文各章只提供活动与主题的工程语境。本附录因此按自足深讲编写——凡指向“第 X 章”处均指主题语境，不意味着正文有逐表逐条的展开。
\end{noteblock}

\Needspace{6\baselineskip}
\section{怎么使用这份术语表（条目模板与受控来源）}

翻开标准的词汇卷，一种常见读法是“用到哪个查哪个”——按字母序找到词条，读完定义就合上。这种读法有一个隐蔽的代价：字母序把概念网剪成了 185 个孤立的点，“fault”排在 f，“failure”也排在 f，两者只隔几行，但它们之间那条“故障显现为失效”的因果箭头在字母序里是看不见的。第 2 章花了整整一章把这张网重新连起来；本附录延续同一套连法，把词条按\textbf{五个概念群}分节，让读者查一个词的时候顺便看见它的邻居。

每个条目固定使用同一个模板：

\begin{itemize}
\item \textbf{中文名（English name, 缩写）}（词条号 ｜ 首讲章）——一句白话，说清它是什么、和邻居差在哪。
\end{itemize}

四个字段各有严格口径：

\begin{enumerate}
\item \textbf{词条号}：形如 1-3.84，表示 ISO 26262-1:2018 第 3 章第 84 个词条。这是每个条目的受控身份，也是与原文核对的唯一入口。本附录定稿前曾用两种独立生成的提取视图交叉定位词条；两种视图来自同一份 PDF，不构成两个独立来源。其中四个词条（3.60、3.153、3.154、3.175）的标题行存在 OCR 噪声，已回看原 PDF 校订。
\item \textbf{首讲章}：指向本书中负责首次完整讲解该术语的章；本附录 C.7.1 给出五群速查，各条目再标明具体去向。术语表只给一句白话；想看完整的直觉铺垫、ISO 定义精读和工程案例，请回到首讲章。个别术语另有概览章与深讲章（如独立性的组织含义在 ch03、技术独立与相关失效在 ch08），条目里会一并标出。
\item \textbf{一句白话}：以原文定义为底、面向工程背景大学生的重述。白话追求“读一遍就懂”，代价是省略了定义中的限定条件——凡是要做合规判断的场合，必须回到原文词条。
\item \textbf{概念群}：一个词条只有一份，但可以出现在多条概念路径上（如 hazard 既是风险群主角、又是时间群时间轴的触发端）。正文按其主路径归群，交叉关系在白话里点出。
\end{enumerate}

机器侧唯一对应物是 \nolinkurl{semantica.chapter_packages.vol2.normative} 中的
\nolinkurl{part1-vocabulary} 与 \nolinkurl{part1-cards} 资产；它们为已迁移词条保存工程释义与来源坐标，
供该包声明的查询和门禁使用。\textbf{注意分母账}：Part 1 的受控词条集共 185 条；
package 当前是 \nolinkurl{partial}、release \nolinkurl{blocked}，必须从当次 manifest/receipt 读取实际覆盖，
不能用局部实例数冒充 185 条全量，也不能把释义资产冒充 ISO 原文。

关于缩写还有一条纪律。本附录在条目里给出的缩写（ASIL、HARA、FTTI、SEooC 等）以原文词条标题或全书首讲章的展开为准；正文遵循“首现三步走”，缩写只在全称出现过之后使用。工程文档需要特别防范\textbf{裸抛缩写}——一份写满 FDTI、MPFDTI、EOTTI 的报告，对没有背景的读者是密文。建议的做法与本书相同：每份文档第一次使用缩写时给出全称与词条号，之后再使用缩写；词条号让读者能沿着本附录回到定义。另外注意有些常用简称（如 FSR、TSR）虽在工程口语中通行，其词条标题是完整短语，缩写地位来自行业惯例——本附录如实标注，不为惯例升格。

还有一条使用纪律值得写在最前面：\textbf{白话与原文的分工是“引路”与“裁决”}。白话负责让你在三秒内想起这个词在概念网里站在哪个位置、和哪个近邻容易混；原文词条负责在争议出现时给出裁决依据。一种需要警惕的错用是拿着白话去做裁决——例如仅凭“安全故障就是无关大局的故障”这句引路语，就把某个未分析的故障归入安全故障，而跳过了原文定义里“不显著提高违背安全目标概率”这个需要论证的判据。每当一个词条要进入 HARA 表格、安全分析报告或供应商协议，请以词条号为键回到正版原文；本附录的词条号可作为检索键，关键使用仍应核对适用版本和原词条。


\Needspace{6\baselineskip}
\section{结构群：从相关项到软件单元}

\begin{noteblock}
首讲 ch02 §2.2。这一群回答“我们给\textbf{什么东西}做安全”：从整车层面的记账单位一路分解到可以单独测试的最小软件颗粒。读这一群的窍门是记住三条轴：\textbf{部分-整体轴}（item\(\rightarrow\)system\(\rightarrow\)component\(\rightarrow\)part/unit）、\textbf{实现轴}（架构\(\rightarrow\)载体技术）、\textbf{边界轴}（谁在相关项里、谁在外面）。
\end{noteblock}

\Needspace{5\baselineskip}
\subsection{分解主干：从记账单位到最小颗粒}

\begin{itemize}
\item \textbf{相关项（item）}（1-3.84 ｜ 首讲 ch02）——在\textbf{整车层面}实现某个功能（或功能的一部分）、被 ISO 26262 当作记账单位的那个系统或系统组合。边界划到哪里，安全责任就算到哪里；HARA、安全目标、安全档案全都以它为单位展开。
\item \textbf{系统（system）}（1-3.163 ｜ 首讲 ch02）——至少把一个传感器、一个控制器和一个执行器关联起来的组件集合。“最小三件套”是它区别于普通组件的判据：少一件，就还不配叫 system。
\item \textbf{元素（element）}（1-3.41 ｜ 首讲 ch02）——系统、组件（硬件或软件）、硬件部件或软件单元的\textbf{统称}。当一句话想同时覆盖各分解层级时就用它；说“软件元素”时指软件组件或软件单元。
\item \textbf{组件（component）}（1-3.21 ｜ 首讲 ch02）——非系统级、逻辑上或技术上可分离的元素，由\textbf{多个硬件部件}，或由\textbf{一个及以上软件单元}组成。硬件分支与软件分支的数量门槛不同，不要误读成两边都必须有多个。
\item \textbf{硬件部件（hardware part）}（1-3.71 ｜ 首讲 ch02）——硬件组件第一层分解出来的部分，如微控制器里的 CPU、一只电阻、一片闪存阵列。
\item \textbf{硬件子部件（hardware subpart）}（1-3.73 ｜ 首讲 ch06）——硬件部件再往下逻辑划分的第二层及更深层，如微控制器 CPU 里的 ALU。
\item \textbf{硬件基本子部件（hardware elementary subpart）}（1-3.72 ｜ 首讲 ch06）——安全分析中考虑的最小硬件颗粒，如带逻辑锥的触发器、一个寄存器。硬件度量的分母就数到这一层为止。
\item \textbf{软件组件（software component）}（1-3.157 ｜ 首讲 ch02）——一个或多个软件单元的集合。定义只有一句话，但它是软件架构分层的基本砖块。
\item \textbf{软件单元（software unit）}（1-3.159 ｜ 首讲 ch02，操作深讲 ch07）——软件架构中\textbf{可以单独测试}的原子级软件组件。“能不能单独测”是判据；Part 6 的单元设计与单元验证都以它为对象。
\item \textbf{嵌入式软件（embedded software）}（1-3.42 ｜ 首讲 ch07）——完成集成、将在处理单元上执行的软件整体。Part 6 末段的“嵌入式软件测试”测的就是这个集成后的整体，而非某个单元。
\end{itemize}

\Needspace{5\baselineskip}
\subsection{架构与实现载体}

\begin{itemize}
\item \textbf{架构（architecture）}（1-3.1 ｜ 首讲 ch05）——相关项或元素的结构表示，能从中识别构造块、块的边界与接口，并包含需求到这些构造块的分配。注意它是“表示”（representation），不是实物本身；原词条说的是需求分配，不是“运行所需过程”。
\item \textbf{安全架构（safety architecture）}（1-3.135 ｜ 首讲 ch05）——为满足安全要求而组织起来的元素集合及其交互。同一个系统的安全架构是其整体架构中承担安全的那张投影。
\item \textbf{电气/电子系统（electrical and/or electronic system, E/E 系统）}（1-3.40 ｜ 首讲 ch02）——由电气或电子元素构成的系统，包括可编程电子元素。ISO 26262 的适用范围就圈在 E/E 系统的功能异常上。
\item \textbf{其他技术（other technology）}（1-3.105 ｜ 首讲 ch02）——E/E 之外的技术，如机械、液压。标准不管它们的开发，但安全机制可以借助它们实现。
\item \textbf{处理单元（processing element）}（1-3.113 ｜ 首讲 appA）——提供数据处理功能的硬件部件，通常含寄存器组、执行单元和控制单元。多核、半导体讨论的基本单位。
\item \textbf{多核（multi-core）}（1-3.95 ｜ 首讲 appA）——含两个及以上可相互独立运行的处理单元的硬件组件。“可独立运行”是它进入 DFA 议题的原因。
\item \textbf{可编程逻辑器件（programmable logic device, PLD）}（1-3.114 ｜ 首讲 appA）——出厂时电路功能未定义、在集成阶段才配置的硬件组件或部件，如 FPGA。
\item \textbf{传感转换器（transducer）}（1-3.172 ｜ 首讲 appA）——把一种能量形式转换为另一种形式的硬件部件，其灵敏度决定输出与被测输入的关系。
\item \textbf{分区（partitioning）}（1-3.106 ｜ 首讲 ch07，深讲 ch08）——为实现某种设计而对功能或元素做的分隔；可用于故障遏制，避免级联失效，是免于干扰的常用手段。
\item \textbf{目标环境（target environment）}（1-3.166 ｜ 首讲 ch07）——特定软件预期运行的环境；对应用软件而言就是那块微控制器加上系统软件。
\end{itemize}

\Needspace{5\baselineskip}
\subsection{边界物种与功能视角}

\begin{itemize}
\item \textbf{候选者（candidate）}（1-3.16 ｜ 首讲 ch10）——定义和使用条件与已发布产品相同或高度一致的相关项/元素——它是“在用证明”论证的主角：想免做部分开发，先证明你是合格的候选者。
\item \textbf{独立于环境的安全要素（safety element out of context, SEooC）}（1-3.138 ｜ 首讲 ch02）——不针对某个具体相关项开发的安全相关元素，如通用微控制器、AUTOSAR 基础软件。它靠“假设的使用条件”参与安全论证，装车时必须逐条核对假设是否成立。
\item \textbf{安全相关元素（safety-related element）}（1-3.144 ｜ 首讲 ch02）——有潜力\textbf{违背或达成}某个安全目标的元素。注意双向措辞：不只是“会闯祸的”，还包括“负责兜底的”。
\item \textbf{安全相关功能（safety-related function）}（1-3.145 ｜ 首讲 ch02）——同上的功能视角版本：有潜力违背或达成安全目标的功能。
\item \textbf{车辆功能（vehicle function）}（1-3.178 ｜ 首讲 ch04）——由一个或多个相关项实现、顾客可观察到的车辆行为，如自适应巡航。功能是顾客视角，相关项是工程记账视角。
\item \textbf{预期功能（intended functionality）}（1-3.83 ｜ 首讲 ch04）——为相关项规定的行为，\textbf{不含安全机制}；规定的行为在整车层面描述。安全机制是为守护预期功能加上的，不算进预期功能本身。
\item \textbf{功能概念（functional concept）}（1-3.66 ｜ 首讲 ch04）——为达成期望行为而对预期功能及其交互的规格说明。它是 HARA 的输入：先说清车要干什么，才能问会怎么出事。
\item \textbf{标定数据（calibration data）}（1-3.15 ｜ 首讲 ch07）——软件构建完成后作为参数值写入的数据，如低怠速目标值、发动机特性曲线。改标定不改代码，但可能改安全行为。
\item \textbf{配置数据（configuration data）}（1-3.22 ｜ 首讲 ch07）——在元素构建阶段赋值、控制构建过程本身的数据，如预处理器变量。与标定数据的分界在于作用时机：一个管“怎么构建”，一个管“构建完怎么调”。
\item \textbf{基线（baseline）}（1-3.10 ｜ 首讲 ch10）——经批准的一组工作产物、相关项或元素的版本，作为后续变更的基准。没有基线，变更管理就没有“变更前”可言。
\end{itemize}

\textbf{路径走查}：拿全书贯穿的 EPS-RC17（电动助力转向候选）把这条路径走一遍。EPS 作为\textbf{相关项}记账——它在整车层面实现“转向助力”这个\textbf{车辆功能}；相关项里的 ECU 凑齐了传感器（扭矩传感器）、控制器（微控制器）、执行器（电机驱动）三件套，够格叫\textbf{系统}；微控制器是一个\textbf{组件}，它内部的 CPU、闪存是\textbf{硬件部件}，CPU 里的 ALU 是\textbf{硬件子部件}；软件侧，扭矩合理性检查是一个可单独测试的\textbf{软件单元}，若干单元组成监控\textbf{软件组件}，全部集成后烧进片子的是\textbf{嵌入式软件}。当供应商拿来一颗按 ASIL D 流程开发、但没针对这台车的通用微控制器时，它是 \textbf{SEooC}——装进相关项边界前，必须逐条核对它的使用假设。一句话可以覆盖上述任何一层的说法，就是\textbf{元素}。分解到哪一层就用哪一层的词，是结构群最基本的语言纪律：把“软件单元”说成“模块”、把“相关项”说成“系统”，在评审桌上都要付出返工的代价。


\Needspace{6\baselineskip}
\section{因果群：故障、错误、失效与依赖失效}

\begin{noteblock}
首讲 ch02 §2.3；\(\lambda\) 五分类深讲 ch06；依赖失效家族深讲 ch08。这一群回答“东西是\textbf{怎么坏}的”：主链 fault\(\rightarrow\)error\(\rightarrow\)failure 每个箭头都写着 can（可能而非必然），旁支按来源、按时长、按后果各切一刀，最后延伸到“多个元素一起坏”的依赖失效家族。
\end{noteblock}

\Needspace{5\baselineskip}
\subsection{主链三词与显现}

\begin{itemize}
\item \textbf{故障（fault）}（1-3.54 ｜ 首讲 ch02）——\textbf{可能}导致元素或相关项失效的异常状态。它是因果链的起点：故障可以潜伏很久而不发作。
\item \textbf{错误（error）}（1-3.46 ｜ 首讲 ch02）——计算、观测或测量出的值/状态，与真实、规定或理论正确的值/状态之间的\textbf{偏差}。它是故障被激活后的中间形态。
\item \textbf{失效（failure）}（1-3.50 ｜ 首讲 ch02）——由于故障显现，元素或相关项的预期行为\textbf{终止}。三个词按“状态\(\rightarrow\)偏差\(\rightarrow\)行为终止”排队，不可混用。
\item \textbf{功能异常表现（malfunctioning behaviour）}（1-3.88 ｜ 首讲 ch02）——相关项相对设计意图的失效或\textbf{非预期行为}。它比失效更宽：行为没终止但偏离了意图，也算。功能安全定义里点名的正是它。
\item \textbf{失效模式（failure mode）}（1-3.51 ｜ 首讲 ch06）——元素或相关项未能提供预期行为的\textbf{方式}。同一个元件可以有多种失效模式（断路、短路、漂移），安全分析按模式逐个记账。
\item \textbf{失效率（failure rate, \(\lambda\)）}（1-3.53 ｜ 首讲 ch06）——失效概率密度除以存活概率，硬件可靠性记账的基本量，通常在使用寿命内假设为常数。
\item \textbf{基础失效率（base failure rate）}（1-3.8 ｜ 首讲 ch06）——给定应用工况下硬件元素的失效率，作为安全分析的\textbf{输入}。它是后续 SPFM/LFM/PMHF 计算链的最上游数字。
\item \textbf{故障模型（fault model）}（1-3.58 ｜ 首讲 appA）——由故障导出失效模式的表示方法，用来评估特定故障的后果，如 stuck-at 模型。
\end{itemize}

\Needspace{5\baselineskip}
\subsection{按来源与时长分}

\begin{itemize}
\item \textbf{随机硬件故障/失效（random hardware fault / random hardware failure）}（1-3.119 / 1-3.118 ｜ 首讲 ch06）——服从概率分布的硬件故障，及其在寿命期内不可预测地发生、但整体遵循概率分布的失效。对付它靠\textbf{度量与机制}（算 \(\lambda\)、加诊断），不靠改流程。
\item \textbf{系统性故障/失效（systematic fault / systematic failure）}（1-3.165 / 1-3.164 ｜ 首讲 ch06，软件治理见 ch07）——以确定性方式与某个原因关联、只能靠\textbf{更改设计或流程}才能消除的故障/失效。软件缺陷全是系统性的；对付它靠开发过程，算不出 \(\lambda\)。
\item \textbf{永久故障（permanent fault）}（1-3.109 ｜ 首讲 ch06）——发生后一直存在、直到移除或修复的故障，如 stuck-at、桥接。
\item \textbf{瞬态故障（transient fault）}（1-3.173 ｜ 首讲 ch06）——发生一次随后消失的故障，如电磁干扰或粒子翻转引起的位翻转。诊断策略对永久/瞬态故障的假设完全不同。
\end{itemize}

\Needspace{5\baselineskip}
\subsection{按后果分：\(\lambda\) 五分类}

\begin{itemize}
\item \textbf{单点故障/单点失效（single-point fault / single-point failure）}（1-3.156 / 1-3.155 ｜ 首讲 ch06）——元素中\textbf{没有任何安全机制覆盖}、自身直接导致违背安全目标的硬件故障，及其导致的失效。注意判据有两半：直接违背 + 无机制覆盖。
\item \textbf{残余故障（residual fault）}（1-3.125 ｜ 首讲 ch06）——元素\textbf{有}安全机制覆盖、但机制没盖住的那一部分随机硬件故障，其自身仍导致违背安全目标。单点与残余的分界只在“有没有机制”，这也是 SPFM 分子里两项并列的原因。
\item \textbf{双点故障/双点失效（dual-point fault / dual-point failure）}（1-3.39 / 1-3.38 ｜ 首讲 ch06）——需要与另一个独立故障组合才导致违背安全目标的单个故障，及两个独立故障组合直接导致的失效。“被监控功能坏了 + 监控它的机制也坏了”是一个用来解释该概念的典型组合，不代表频率排名。
\item \textbf{多点故障/多点失效（multiple-point fault / multiple-point failure）}（1-3.97 / 1-3.96 ｜ 首讲 ch06）——推广到 n 个独立故障组合的版本。若未被检测、未被感知，多点故障就在系统里潜伏着等队友。
\item \textbf{潜伏故障（latent fault）}（1-3.85 ｜ 首讲 ch06）——在多点故障检测时间间隔内，既没被安全机制检出、也没被驾驶员感知的多点故障。它是 LFM 的分子对象：机制自己坏了而未被发现，是一个典型的教学例。
\item \textbf{被检出故障（detected fault）}（1-3.31 ｜ 首讲 ch06）——在规定时间内被安全机制检出的故障。“规定时间”通常就是时间群里的故障检测时间间隔。
\item \textbf{被感知故障（perceived fault）}（1-3.108 ｜ 首讲 ch06）——通过整车层面的行为偏差被\textbf{间接}察觉的故障，如方向盘变沉。项目若要用驾驶员感知支持故障分类，应说明在所主张的场景中行为偏差可被感知、代表性人员能识别并作出响应，且感知和响应时间满足相关 MPFDTI、FTTI 或其他已规定时间窗口。它也可影响诊断、HMI、警告与服务策略的假设；词条本身不使这些假设自动成立。
\item \textbf{安全故障（safe fault）}（1-3.130 ｜ 首讲 ch06）——发生后不会显著提高违背安全目标概率的故障。名字容易误导：它不是“安全机制的故障”，而是“无关大局的故障”。
\item \textbf{故障容忍（fault tolerance）}（1-3.60 ｜ 首讲 ch05）——在一个或多个规定故障存在的情况下，仍能交付规定功能的能力。冗余架构买的就是这个能力。
\end{itemize}

\Needspace{5\baselineskip}
\subsection{依赖失效家族}

\begin{itemize}
\item \textbf{相关失效（dependent failures）}（1-3.29 ｜ 首讲 ch08）——统计上不独立的多个失效：组合发生的概率不等于各自概率的简单乘积。冗余的算术全建立在独立假设上，相关失效就是拆这个台的。
\item \textbf{独立失效（independent failures）}（1-3.79 ｜ 首讲 ch08）——同时或相继发生的概率恰好等于各自无条件概率乘积的失效。它是相关失效的对照组，也是概率计算的前提假设。
\item \textbf{级联失效（cascading failure）}（1-3.17 ｜ 首讲 ch08）——某个相关项中的元素因内外部根因失效，进而\textbf{引起}同一相关项或不同相关项中的另一个或多个元素失效。A 推倒 B，关键是有传播方向，不限于同一相关项内部。
\item \textbf{共因失效（common cause failure, CCF）}（1-3.18 ｜ 首讲 ch08）——两个及以上元素\textbf{直接}由同一个特定事件或根因导致的失效。同一电源掉电导致主备双双失效是标准例子；与级联的区别在于“同一根因直接打到多个元素”而非逐个传导。
\item \textbf{共模失效（common mode failure）}（1-3.19 ｜ 首讲 ch08）——CCF 的特例：多个元素以\textbf{相同方式}失效；何种相似程度足以归为共模，要结合语境判断，并不要求失效表现完全相同。若两个传感器因同一批次缺陷这一共同根因而以相近方式漂移，可构成共模失效。
\item \textbf{耦合因素（coupling factors）}（1-3.26 ｜ 首讲 ch08）——导致多个元素失效相关联的共同特征或关系，如同一供电、同一位置、同一开发团队。DFA 可以按已声明的分析边界和耦合因素分类，系统搜索可信的相关失效路径；清单不自动代表搜索穷尽。
\item \textbf{相关失效引发点（dependent failure initiator, DFI）}（1-3.30 ｜ 首讲 ch08）——通过耦合因素让多个元素失效的\textbf{单一根因}。DFA 应在声明边界内识别可信的 DFI 候选及所影响元素，并记录消除、控制或接受理由；不能从“已列清单”推出所有发起者已被找到或挡住。
\item \textbf{独立性（independence）}（1-3.78 ｜ 组织含义首讲 ch03，技术含义首讲 ch08）——两重含义：元素之间不存在会违背安全要求的相关失效；或组织之间不存在利益冲突。ASIL 分解的前提、确认措施的资格，用的都是这个词。
\item \textbf{免于干扰（freedom from interference, FFI）}（1-3.65 ｜ 首讲 ch08）——两个及以上元素之间不存在会导致违背安全要求的\textbf{级联失效}。软件分区要证明的就是它：QM 分区不得拖垮 ASIL 分区。正文以“互不干扰”的叙事语汇出现，合规判据以本词条为准。
\item \textbf{安全异常（safety anomaly）}（1-3.134 ｜ 首讲 ch09）——偏离预期且可能导致伤害的状况。评审、测试、现场运行中发现的“不对劲”都先记为安全异常，再决定处置。
\end{itemize}

\Needspace{5\baselineskip}
\subsection{覆盖率与观测}

\begin{itemize}
\item \textbf{诊断覆盖率（diagnostic coverage, DC）}（1-3.33 ｜ 首讲 ch06）——硬件元素失效率中被安全机制\textbf{检出或控制}的百分比。它是把“装了诊断”翻译成数字的桥梁。
\item \textbf{失效模式覆盖率（failure mode coverage）}（1-3.52 ｜ 首讲 ch06）——某失效模式的失效率中被安全机制检出或控制的比例。DC 按元素记总账，它按失效模式记明细账。
\item \textbf{硬件架构度量（hardware architectural metrics）}（1-3.70 ｜ 首讲 ch06）——评价硬件架构对安全的有效性的度量，即 SPFM 与 LFM 的统称。
\item \textbf{诊断点（diagnostic points）}（1-3.34 ｜ 首讲 appA）——观察到故障被\textbf{检测或纠正}的元素输出信号。
\item \textbf{观测点（observation points）}（1-3.101 ｜ 首讲 appA）——观察到故障\textbf{潜在影响}的元素输出信号，如存储器的输出。诊断点看“抓没抓住”，观测点看“有没有影响”。
\end{itemize}

\textbf{路径走查}：EPS 扭矩传感器的信号线发生桥接短路——这是一个\textbf{永久}的\textbf{随机硬件故障}；它让 ADC 读出的扭矩值偏离真实值，这是\textbf{错误}；错误穿过滤波与控制律，最终让助力电机输出非预期扭矩、“按需助力”这一预期行为终止，这是\textbf{失效}，其\textbf{失效模式}是“助力扭矩超出请求值”。现在按后果分类记账：如果这条信号链没有任何安全机制覆盖，它就是\textbf{单点故障}；如果装了合理性检查但覆盖率只有 90\%，漏网的 10\% 是\textbf{残余故障}；合理性检查电路自己坏了而无人察觉，则是一个正在潜伏的\textbf{多点故障}——若在\textbf{多点故障检测时间间隔}内既没被检出也没被驾驶员感知，就归入\textbf{潜伏故障}。是否已被感知不能只凭故障名称假定，要用所主张场景中的可感知性、驾驶员响应和时间窗口证据支持。再横向看一眼：主备两路传感器共用同一 5 V 供电，供电跌落让两路同时漂移——这不是两个独立失效的巧合，而是以供电为\textbf{耦合因素}、由同一个\textbf{相关失效引发点}触发的\textbf{共因失效}。因果群的词汇帮助把这段叙述拆成可分析、可追溯的对象，具体归类仍取决于项目边界、故障模型与证据。


\Needspace{6\baselineskip}
\section{风险群：危害事件、S/E/C 与完整性等级}

\begin{noteblock}
首讲 ch02 §2.5；操作深讲 ch04。这一群回答“坏了\textbf{有多危险}”：从伤害一路走到 ASIL。路线上最关键的一步是“危害 \(\times\) 运行场景 = 危害事件”——分级的对象从来不是部件，而是整车层面的危害事件。
\end{noteblock}

\Needspace{5\baselineskip}
\subsection{从危害到危害事件}

\begin{itemize}
\item \textbf{安全（safety）}（1-3.132 ｜ 首讲 ch02）——不存在不合理的风险。整本标准的地基只有这一句：安全不是零风险。
\item \textbf{功能安全（functional safety）}（1-3.67 ｜ 首讲 ch02）——不存在由 E/E 系统\textbf{功能异常表现}引起的危害所导致的不合理风险。三个限定词划定了整本书的边界：对象是 E/E 系统的功能异常及其危害，目标是不存在不合理风险，不是宣称把风险消除为零。
\item \textbf{伤害（harm）}（1-3.74 ｜ 首讲 ch04）——对人的身体损伤或健康损害。注意对象只有人：撞坏护栏本身不算 harm。
\item \textbf{危害（hazard）}（1-3.75 ｜ 首讲 ch04）——由相关项功能异常表现引起的\textbf{潜在}伤害来源。它是“可能性”而非“事件”：非预期的转向助力是危害，撞车才是事故。
\item \textbf{非功能性危害（non-functional hazard）}（1-3.100 ｜ 首讲 ch04）——由 E/E 功能异常之外的因素（如电击、火灾、毒性）引起的危害。ISO 26262 明确不管它们，划界时要说得出去哪儿管。
\item \textbf{运行场景（operational situation）}（1-3.104 ｜ 首讲 ch04）——车辆生命期内可能出现的场景，如高速行驶、坡道停车、维修。它是 E 参数的评估对象。
\item \textbf{运行模式（operating mode）}（1-3.102 ｜ 首讲 ch04）——相关项或元素因使用而处的功能状态，如系统关闭、正常运行、降级运行。
\item \textbf{车辆运行状态（vehicle operating state）}（1-3.179 ｜ 首讲 ch04）——运行模式与运行场景的组合。场景在车外、模式在车内，两者合起来才描述“车此刻的处境”。
\item \textbf{危害事件（hazardous event）}（1-3.77 ｜ 首讲 ch04）——危害与运行场景的\textbf{组合}。同一个危害配上不同场景，严重度和可控性都可能不同——这是 HARA 按事件而非按危害逐行打分的原因。
\item \textbf{危害分析与风险评估（hazard analysis and risk assessment, HARA）}（1-3.76 ｜ 首讲 ch04）——识别并分级相关项的危害事件、进而制定安全目标及其 ASIL 的方法。它是概念阶段的发动机。
\end{itemize}

\Needspace{5\baselineskip}
\subsection{三参数与风险}

\begin{itemize}
\item \textbf{严重度（severity, S）}（1-3.154 ｜ 首讲 ch04）——潜在危害事件中一人或多人可能受到伤害程度的\textbf{估计}。HARA 中取 S0–S3。
\item \textbf{暴露（exposure, E）}（1-3.48 ｜ 首讲 ch04）——处于某个运行场景中的状态，该场景若与被分析的失效模式同时出现则可能有危害。E 参数评估的是\textbf{场景出现的可能性}（正文因此把这项判断称作“暴露概率”），不是失效的可能性——这是本附录要重点辨析的区别。
\item \textbf{可控性（controllability, C）}（1-3.25 ｜ 首讲 ch04）——相关人员通过及时反应（可能借助外部措施）避免特定伤害或损失的能力。相关人员可包括驾驶员、乘员或车辆外部附近人员；在具体 HARA 场景中应识别真正涉事的人，不能把定义直接缩成“典型驾驶员”。
\item \textbf{专家骑手（expert rider）}（1-3.47 ｜ 首讲 appB）——摩托车领域中有能力基于实车操作评估可控性分级的角色。Part 12 用它校准摩托车的 C 参数。
\item \textbf{风险（risk）}（1-3.128 ｜ 首讲 ch04）——伤害发生\textbf{概率}与该伤害\textbf{严重度}的组合。两个因子缺一不可。
\item \textbf{不合理风险（unreasonable risk）}（1-3.176 ｜ 首讲 ch02）——依据当时有效的社会道德观念，在特定情境下被判定为不可接受的风险。注意判据是社会共识而非工程计算——它解释了为什么“安全”没有全球统一的数值阈值。
\item \textbf{残余风险（residual risk）}（1-3.126 ｜ 首讲 ch04）——部署安全措施\textbf{之后}剩下的风险。安全工程的目标是把残余风险压到不合理线以下，而不是压到零。
\item \textbf{合理可预见（reasonably foreseeable）}（1-3.120 ｜ 首讲 ch04）——技术上可能、且发生率可信或可测的情形。可预见的误用要管，科幻式的滥用不用管，界线就画在这个词上。
\item \textbf{外部措施（external measure）}（1-3.49 ｜ 首讲 ch04）——独立于相关项之外、能降低或缓解相关项风险的措施，如道路护栏。HARA 打分时它能救 C 参数一命，但不算相关项自己的安全机制。
\end{itemize}

\Needspace{5\baselineskip}
\subsection{完整性等级}

\begin{itemize}
\item \textbf{汽车安全完整性等级（automotive safety integrity level, ASIL）}（1-3.6 ｜ 首讲 ch02，操作深讲 ch04）——四个等级（A 最低、D 最高）之一，用来规定相关项或元素为避免不合理风险而需适用的 ISO 26262 要求与安全措施。它把危害事件的风险分级连接到工程严谨度；原词条说的是 unreasonable risk，不是另造一个“不合理残余风险”术语。
\item \textbf{ASIL 能力（ASIL capability）}（1-3.2 ｜ 首讲 ch10）——相关项或元素满足\textbf{假设的}、按给定 ASIL 分配的安全要求的能力。SEooC 的说法：“我按 D 的流程开发过”说的是能力，不是这次装车就自动合格。
\item \textbf{ASIL 分解（ASIL decomposition）}（1-3.3 ｜ 首讲 ch08）——把冗余的安全要求分配到\textbf{充分独立}的多个元素上、服务于同一安全目标，从而降低各元素 ASIL 的做法。注意三个不可省略的限定：冗余、充分独立、同一安全目标——独立性证据不到位，分解就是自欺。
\item \textbf{摩托车安全完整性等级（motorcycle safety integrity level, MSIL）}（1-3.94 ｜ 首讲 appB）——摩托车领域的四级风险降低等级，需按 Part 12 规则\textbf{换算}为 ASIL 后再套用各部要求。同名不同域，不可直接混用。
\item \textbf{质量管理（quality management, QM）}（1-3.117 ｜ 首讲 ch04）——组织为质量而进行的协调指挥与控制活动。Part 1 的注记只明确两点：\textbf{QM 不是 ASIL}，但可以在 HARA 中指定；至于某个 HARA 结果后续适用哪些活动，应回到 Part 3 及项目适用性判断，不能从词条本身推出“无需 ISO 26262 附加措施”这一无条件结论。
\end{itemize}

\textbf{路径走查}（教学假设，非任何项目结论）：继续 EPS 的故事。“非预期的转向助力”是一个\textbf{危害}——潜在的伤害来源，还不是危害事件。把它分别与“高速过弯”和“停车场低速挪车”等\textbf{运行场景}组合，才得到需要在 \textbf{HARA} 表中分行评估的\textbf{危害事件}。为演示查表，假设项目已用指定车速、道路曲率、碰撞对象与伤情谱支持 \textbf{S3}，用目标车队的运行场景数据支持 \textbf{E4}，并用代表性驾驶员、故障施加时刻、可用纠正动作与反应时间窗口的评估记录支持 \textbf{C3}。在这组给定教学输入下，S3+E4+C3 查表得 \textbf{ASIL D}；本附录不提供上述项目证据，因而不为现实 EPS 场景作等级断言。这段走查要区分的是：E 评场景出现的可能性而不是故障概率，C 评场景中相关人员避免特定伤害的能力，S 估计伤害程度而不是维修费。相似安全目标合并时，取其所关联危害事件中的最高 ASIL。若某行在 HARA 中被指定为 \textbf{QM}，表示该结果不是 ASIL；后续活动应按 Part 3、项目适用性和组织质量管理体系判断，不能把 QM 简化成“没有其他义务”。


\Needspace{6\baselineskip}
\section{需求群：安全目标、需求与工作产物}

\begin{noteblock}
概览 ch02；操作深讲 ch04/ch05，档案与确认深讲 ch03。这一群回答“我们\textbf{答应做到什么}、拿什么证明”：需求沿“安全目标\(\rightarrow\)功能安全要求\(\rightarrow\)技术安全要求\(\rightarrow\)硬件/软件实现”逐层落地；措施与机制负责兑现承诺；工作产物与各种评审、审核构成证据链。词汇表里没有 HSR 和 SSR——这条诚实边界由本附录维持，此处不再造词。
\end{noteblock}

\Needspace{5\baselineskip}
\subsection{需求层级}

\begin{itemize}
\item \textbf{安全目标（safety goal, SG）}（1-3.139 ｜ 首讲 ch04）——HARA 得出的\textbf{整车层面最高层}安全要求。一个安全目标可以对应多个危害事件，其 ASIL 取所涉事件中最高者。
\item \textbf{功能安全概念（functional safety concept, FSC）}（1-3.68 ｜ 首讲 ch04）——功能安全要求的规格说明及其到相关项架构元素的分配。它回答“用什么\textbf{策略}兑现安全目标”。
\item \textbf{功能安全要求（functional safety requirement, FSR）}（1-3.69 ｜ 首讲 ch04）——\textbf{与实现无关}的安全行为或安全措施的规格说明。判据就在“implementation-independent”这一个词上：说了“用哪种传感器”就已经不是 FSR。
\item \textbf{技术安全概念（technical safety concept, TSC）}（1-3.167 ｜ 首讲 ch05）——技术安全要求的规格说明及其到系统元素的分配，附带实现所需的信息。
\item \textbf{技术安全要求（technical safety requirement, TSR）}（1-3.168 ｜ 首讲 ch05）——为\textbf{实现}关联的功能安全要求而导出的要求。FSR 说“要防非预期助力”，TSR 说“扭矩指令超阈值 x ms 内切断”——从策略到技术指标的翻译。
\item \textbf{继承（inheritance）}（1-3.81 ｜ 首讲 ch05）——开发过程中把需求属性\textbf{原封不动}传递到下一细化层级。ASIL 沿需求链往下传，默认走的就是继承，除非启动 ASIL 分解。
\end{itemize}

\Needspace{5\baselineskip}
\subsection{措施与机制}

\begin{itemize}
\item \textbf{安全措施（safety measure）}（1-3.141 ｜ 首讲 ch02）——避免或控制系统性失效、检测或控制随机硬件失效或缓解其危害影响的\textbf{活动或技术方案}。范围很宽：编码规范、评审、看门狗都算。
\item \textbf{安全机制（safety mechanism）}（1-3.142 ｜ 首讲 ch02）——由 E/E 功能/元素或其他技术实现的\textbf{技术方案}，用于检测并缓解或容忍故障、控制或避免失效，以保持预期功能或达到安全状态。机制是措施的真子集：只有做进产品里的技术方案才叫机制，流程活动不算。
\item \textbf{专用措施（dedicated measure）}（1-3.27 ｜ 首讲 ch06）——确保安全目标违背概率评估中所声称失效率成立的措施，如设计特性、专项样本测试、老化筛选。声称 \(\lambda\) 有多低，就要有对应的专用措施撑住。
\item \textbf{冗余（redundancy）}（1-3.122 ｜ 首讲 ch08）——在足以完成功能的手段\textbf{之外}再增加的手段，如双通道采集。冗余是分解的物质基础，但冗余本身不自动等于独立。
\item \textbf{多样性（diversity）}（1-3.37 ｜ 首讲 ch08）——用\textbf{不同的}方案满足同一需求、以谋求独立性的做法。词条注记提醒：多样性不保证不发生共因失效，只是降低其可能。
\item \textbf{稳健设计（robust design）}（1-3.129 ｜ 首讲 ch09）——在无效输入或恶劣环境条件下仍能正确工作的设计；对软件即对无效输入的耐受，对硬件即对环境应力的耐受。
\item \textbf{故障注入（fault injection）}（1-3.57 ｜ 首讲 ch07）——向元素中插入故障、错误或失效，观察其反应以评估故障影响的\textbf{方法}。它是方法表里的常客：单元、集成、整车各层都有它的行。
\item \textbf{基于模型的开发（model-based development）}（1-3.90 ｜ 首讲 ch07）——用模型描述待开发元素行为或属性的开发方式。模型可以生成代码，也因此把工具链的置信度问题（Part 8）带进了软件章。
\end{itemize}

\Needspace{5\baselineskip}
\subsection{工作产物与安全档案}

\begin{itemize}
\item \textbf{工作产物（work product）}（1-3.185 ｜ 首讲 ch03）——由 ISO 26262 一个或多个相关要求产生的\textbf{文档化结果}；它可以是一份完整文档，也可以由多份文档共同承载。是否充分证明活动已执行，要结合相应要求、过程记录和评审判断，不能把“没文档等于没做”当成词条定义的一部分。
\item \textbf{安全档案（safety case）}（1-3.136 ｜ 首讲 ch03）——论证相关项或元素已达成功能安全的\textbf{论证}（argument），由工作产物汇编的证据支撑。注意词性：它是论证结构，不是文件夹——把工作产物堆在一起不构成安全档案。
\item \textbf{安全计划（safety plan）}（1-3.143 ｜ 首讲 ch03）——管理和引导项目安全活动执行的计划，含日期、里程碑、任务、交付物与责任人。
\item \textbf{安全活动（safety activity）}（1-3.133 ｜ 首讲 ch03）——在安全生命周期一个或多个阶段/子阶段中执行的活动。
\item \textbf{安全生命周期（lifecycle）}（1-3.86 ｜ 首讲 ch03）——相关项从概念到退役的全部阶段。
\item \textbf{阶段（phase）}（1-3.110 ｜ 首讲 ch03）——安全生命周期中由 Part 3–7 规定的阶段，如概念阶段、系统层开发。
\item \textbf{子阶段（sub-phase）}（1-3.161 ｜ 首讲 ch03）——阶段在某个条款中规定的细分，如 HARA 是概念阶段的一个子阶段。
\item \textbf{安全经理（safety manager）}（1-3.140 ｜ 首讲 ch03）——负责监督并确保达成功能安全所需活动得以执行的人或组织。
\item \textbf{安全文化（safety culture）}（1-3.137 ｜ 首讲 ch03）——组织中持久的价值观、态度、动机与知识，使安全在决策和行为中\textbf{优先于竞争性目标}。它决定了流程条文在压力下是否还算数。
\item \textbf{管理体系（management system）}（1-3.87 ｜ 首讲 ch03）——组织为达成目标所用的方针、程序与过程的总称。
\end{itemize}

\Needspace{5\baselineskip}
\subsection{验证、确认与各种“看一遍”}

\begin{itemize}
\item \textbf{验证（verification）}（1-3.180 ｜ 首讲 ch05，框架深讲 ch10）——判定被检对象\textbf{是否满足其规定要求}。问的是“做得对不对”（against spec）。
\item \textbf{验证评审（verification review）}（1-3.181 ｜ 首讲 ch10）——确保开发活动结果满足项目要求或技术要求的验证活动。
\item \textbf{安全确认（safety validation）}（1-3.148 ｜ 首讲 ch05）——基于检查和试验，对安全目标\textbf{是否充分，以及是否以足够完整性得到实现}给出确信。它与验证的对象和判据不同：这里问的是安全目标是否适当、是否已经达成，不能简单说成“比验证高一层”。
\item \textbf{测试（testing）}（1-3.169 ｜ 首讲 ch07）——为验证满足要求、检测异常、建立信心而对相关项或元素进行规划、准备、运行或施加激励的过程。
\item \textbf{评审（review）}（1-3.127 ｜ 首讲 ch03）——按评审目的，对工作产物是否达成其预期目标进行的检查。
\item \textbf{审查（inspection）}（1-3.82 ｜ 方法表见附录 D）——\textbf{遵循正式规程}检查工作产物以检出安全异常。与走查同为验证手段，区别在正式程度：审查有规程、角色和进出条件。
\item \textbf{走查（walk-through）}（1-3.182 ｜ 首讲 ch08）——系统地过一遍工作产物以检出安全异常，形式比审查轻。方法表里两者是备选关系（Part 6 Table 7 的 1a/1c）。
\item \textbf{确认措施（confirmation measure）}（1-3.23 ｜ 首讲 ch03）——针对功能安全的确认评审、审核或评估的统称。执行人的独立性等级（I0–I3）随 ASIL 升高，首讲章有完整矩阵。
\item \textbf{确认评审（confirmation review）}（1-3.24 ｜ 首讲 ch03）——确认某工作产物为达成功能安全提供了\textbf{充分且有说服力的证据}。对象是工作产物。
\item \textbf{功能安全审核（audit）}（1-3.5 ｜ 首讲 ch03）——对\textbf{已实施过程}是否符合过程目标的检查。对象是过程。
\item \textbf{功能安全评估（assessment）}（1-3.4 ｜ 首讲 ch03）——对相关项或元素的特性是否达成 ISO 26262 目标的检查。对象是达成的功能安全本身；评审、审核、评估三者对象不同，不可互相顶替。
\item \textbf{回归策略（regression strategy）}（1-3.123 ｜ 首讲 ch10）——验证变更没有影响未变更且已验证部分的策略。改一行代码要不要重测全部，答案写在这里。
\end{itemize}

\Needspace{5\baselineskip}
\subsection{表示法与结构覆盖率}

\begin{itemize}
\item \textbf{形式化/半形式化/非形式化表示法（formal / semi-formal / informal notation）}（1-3.63 / 1-3.149 / 1-3.80 ｜ 首讲 ch07）——按\textbf{语法和语义的定义完备程度}递减的三档描述技术：形式化两者皆完备（如模型检查器输入语言）；半形式化语法完备、语义可不完备（如 UML 子集）；非形式化连语法都不完备（如自然语言配图）。方法表按 ASIL 调节三档的推荐度。
\item \textbf{形式化/半形式化验证（formal / semi-formal verification）}（1-3.64 / 1-3.150 ｜ 首讲 ch07）——形式化验证使用形式化表示法，证明相关项或元素相对于其功能或属性规格的正确性；半形式化验证则基于半形式化表示法给出的描述实施验证，例如用半形式化模型生成测试向量并检查系统行为。两者首先区分的是表示法与方法依据，不能脱离具体方法和验证对象，把差异一概简化成“证明”与“检查”的证据等级。
\item \textbf{语句覆盖率（statement coverage）}（1-3.160 ｜ 首讲 ch07）——测试中被执行的软件语句百分比。三种结构覆盖率中最弱的一档。
\item \textbf{分支覆盖率（branch coverage）}（1-3.13 ｜ 首讲 ch07）——测试中被执行的控制流分支百分比；100\% 分支覆盖蕴含 100\% 语句覆盖，反之不然。
\item \textbf{修正条件/判定覆盖率（modified condition/decision coverage, MC/DC）}（1-3.92 ｜ 首讲 ch07）——独立影响判定结果的单个条件取值被演练到的百分比，三档中最强，Part 6 Table 9 只对 ASIL D 给 ++。
\end{itemize}

\Needspace{5\baselineskip}
\subsection{供应链、工具与既有信任}

\begin{itemize}
\item \textbf{开发接口协议（development interface agreement, DIA）}（1-3.32 ｜ 首讲 ch10）——客户与供应商之间的协议，规定分布式开发中双方各自负责执行的活动、评审的证据与交换的工作产物。安全责任在供应链上的“分界合同”。
\item \textbf{供货协议（supply agreement）}（1-3.162 ｜ 首讲 ch10）——客户与供应商间关于活动、证据或工作产物的责任约定；量产供货语境下与 DIA 相邻但不相同。
\item \textbf{分布式开发（distributed development）}（1-3.36 ｜ 首讲 ch10）——相关项或元素的开发责任在客户与供应商之间划分的开发形态。DIA 因它而生。
\item \textbf{软件工具（software tool）}（1-3.158 ｜ 首讲 ch10）——用于相关项或元素\textbf{开发过程}的计算机程序。编译器、代码生成器、静态分析器都是；它们不装车，但可能把错误注入装车的东西——这是工具置信度（TCL）议题的起点。
\item \textbf{在用证明论证/在用证明信用（proven in use argument / proven in use credit, PiU）}（1-3.115 / 1-3.116 ｜ 提名 ch09，深讲 ch10）——前者是基于候选者现场数据分析所得的证据，用来说明任何可能损害安全目标的候选者失效概率满足\textbf{适用 ASIL 的要求}；后者是在此论证成立后，用它及相应工作产物替代给定生命周期子阶段。不能把适用 ASIL 的整套要求缩成一个脱离语境的“低于某目标值”。
\item \textbf{现场数据（field data）}（1-3.62 ｜ 首讲 ch09）——相关项或元素使用中获得的数据，含累计运行小时、全部失效与在役安全异常。在用证明的原料。
\item \textbf{久经考验（well-trusted）}（1-3.184 ｜ 首讲 ch10）——曾在可比应用中使用且无已知安全异常，如久经考验的设计原则、工具、软件模块。口语的“用过很多年”必须落到“可比应用 + 无已知异常”两个判据上。
\item \textbf{安全相关特殊特性（safety-related special characteristic）}（1-3.147 ｜ 首讲 ch09）——相关项、元素或其\textbf{生产过程}的特性，其合理可预见的偏差可能影响、促成或引起功能安全的潜在降低。它不只覆盖会直接违背安全目标的偏差，也是安全要求从开发端伸进工厂的重要接口。
\end{itemize}

\textbf{路径走查}：需求群是一条自上而下的翻译链，配一条自下而上的证据链。往下走：HARA 给出\textbf{安全目标}“防止非预期助力（ASIL D）”；\textbf{功能安全概念}给出策略——“监控扭矩指令，超限即切断”，写成与实现无关的\textbf{功能安全要求}并分配到架构元素；系统层把策略翻译成\textbf{技术安全要求}——“指令超阈值持续 x ms 内切断电机供电”，汇入\textbf{技术安全概念}；ASIL 沿链\textbf{继承}，除非启动分解。兑现承诺的是\textbf{安全机制}（扭矩合理性检查这个做进产品的技术方案）和更广义的\textbf{安全措施}（包括编码规范、评审这些流程活动）。往上走：每一步活动留下\textbf{工作产物}；\textbf{验证}逐层确认“做得对不对”，\textbf{安全确认}在整车上确认“目标本身对不对、真达到没有”；\textbf{确认评审、审核、评估}三种\textbf{确认措施}由具备独立性的人分别检查工作产物、过程与达成的功能安全；相关证据最终编入\textbf{安全档案}——注意它是论证不是文件夹。链条两端接上，“承诺—兑现—证明”才闭环。


\Needspace{6\baselineskip}
\section{时间群：容忍、检测、处理与运行时间}

\begin{noteblock}
首讲 ch02 §2.4；预算与分配深讲 ch05；现场时序深讲 ch09。这一群回答“来得及吗”。所有词条挂在同一条事件时间轴上：\textbf{故障发生 \(\rightarrow\) 检测 \(\rightarrow\) 反应 \(\rightarrow\) 安全状态（或应急运行 \(\rightarrow\) 安全状态）}。读这一群务必分清两类归属：FTTI 是\textbf{相关项}的属性（留给你多少时间），FDTI/FRTI/FHTI 是\textbf{安全机制}的属性（你实际用掉多少时间）。Part 1 的词条注记用二者的比较解释何谓及时覆盖；具体项目仍应在安全要求和时序论证中规定并验证预算，不能把术语注记冒充独立的 shall。
\end{noteblock}

\begin{itemize}
\item \textbf{故障容忍时间间隔（fault tolerant time interval, FTTI）}（1-3.61 ｜ 首讲 ch02，预算深讲 ch05）——在安全机制未激活的情况下，从故障在相关项中发生到\textbf{可能}出现危害事件的最短时间。它是留给所有防线的总预算，属于相关项。
\item \textbf{故障检测时间间隔（fault detection time interval, FDTI）}（1-3.55 ｜ 首讲 ch02）——从故障发生到被检出的时间。
\item \textbf{故障反应时间间隔（fault reaction time interval, FRTI）}（1-3.59 ｜ 首讲 ch02）——从故障被检出到达到安全状态或进入应急运行的时间。
\item \textbf{故障处理时间间隔（fault handling time interval, FHTI）}（1-3.56 ｜ 概念首讲 ch02，预算深讲 ch05）——FDTI 与 FRTI 之\textbf{和}，即安全机制从故障发生到完成反应的用时。Part 1 对 FDTI 的 Note 4 说明：若 FDTI + FRTI 小于相关 FTTI，则相应安全机制及时覆盖该故障；这是资料性的术语说明，项目义务仍须追到适用的规范性条款和安全要求。
\item \textbf{诊断测试时间间隔（diagnostic test time interval, DTTI）}（1-3.35 ｜ 首讲 ch06）——安全机制两次在线诊断测试执行之间的时间，且包含一次在线诊断测试的执行时长。它描述诊断执行间隔；允许的周期或上限要由具体诊断设计、FDTI/FTTI 分配和安全要求论证，不能从 DTTI 词条单独推出。
\item \textbf{多点故障检测时间间隔（multiple-point fault detection time interval, MPFDTI）}（1-3.98 ｜ 首讲 ch06）——在多点故障可能凑成多点失效\textbf{之前}将其检出的时间窗。潜伏故障的定义就以它为限。
\item \textbf{应急运行（emergency operation）}（1-3.43 ｜ 首讲 ch05）——故障反应之后、过渡到安全状态之前，为维持安全而设的相关项运行模式。直接断电不安全的系统（如行驶中的转向助力）需要这段“降级续命”。
\item \textbf{应急运行时间间隔（emergency operation time interval, EOTI）}（1-3.44 ｜ 概览 ch02，预算深讲 ch05）——应急运行\textbf{实际维持}的时间段。
\item \textbf{应急运行容忍时间间隔（emergency operation tolerance time interval, EOTTI）}（1-3.45 ｜ 概览 ch02，预算深讲 ch05）——应急运行在\textbf{不产生不合理风险}的前提下可以维持的规定时长。EOTI 必须收敛在 EOTTI 之内；一对“实际用时 vs 容忍上限”的关系，与 FHTI/FTTI 同构。
\item \textbf{最大修复时间间隔（maximum time to repair time interval）}（1-3.89 ｜ 首讲 ch09）——安全状态可以维持的规定时长；当安全状态需要靠维修才能退出时，这个窗口决定了“多久内必须修好”。
\item \textbf{安全状态（safe state）}（1-3.131 ｜ 概览 ch02，系统策略深讲 ch05）——发生失效时相关项\textbf{不存在不合理风险}的运行模式。时间轴的终点站。注意它是运行模式而非“关机”：对某些系统，关机恰恰不安全。
\item \textbf{降级（degradation）}（1-3.28 ｜ 首讲 ch09）——相关项或元素功能或性能降低的状态或向该状态的过渡。限速跛行回家是典型降级。
\item \textbf{警告与降级策略（warning and degradation strategy）}（1-3.183 ｜ 首讲 ch09）——规定如何提醒驾驶员功能可能降低、以及如何提供降级功能以到达安全状态的规格说明。时间轴上“人”这条支线的剧本。
\item \textbf{运行时间（operating time）}（1-3.103 ｜ 首讲 ch06）——相关项或元素处于工作状态的累计时间，含降级模式。PMHF 按每小时记账，分母就是它。
\end{itemize}

\textbf{路径走查}：给时间轴配一次算术（数值为教学示例，非任何产品结论）。假设项目在危害分析基础上，为 EPS 相应安全目标规定：非预期助力持续超过 100 ms 才可能酿成危害事件——\textbf{FTTI = 100 ms}，这是相关项给所有防线的总预算。安全机制这边：扭矩合理性检查每 10 ms 跑一轮（\textbf{诊断测试时间间隔} 10 ms），最坏情况下故障发生后 20 ms 被检出（\textbf{FDTI} = 20 ms），再用 30 ms 切断电机供电、进入无助力状态（\textbf{FRTI} = 30 ms），合计 \textbf{FHTI} = 50 ms < 100 ms，不等式成立，且留了 50 ms 教学裕量。若这台车行驶中直接失去助力并不安全，就不能一步跳到安全状态：先进入限制助力的\textbf{应急运行}，实际维持的 \textbf{EOTI} 必须收敛在\textbf{EOTTI}——“降级续命不产生不合理风险”的容忍上限——之内，最终停靠\textbf{安全状态}。这次走查要重点辨析三件事：FTTI 属于相关项，而非某个机制的独立指标；FHTI 由从故障发生起算的 FDTI 与从检出起算的 FRTI 组成；安全状态是具体系统在特定场景下不存在不合理风险的运行模式，不能被概括成“关机”。


\Needspace{6\baselineskip}
\section{本体化实践：一份词表、多条概念路径}

本附录由两种互补载体共同维护：Markdown 正文承担五群讲解、字母索引与首讲回指；
Semantica normative package 的 \nolinkurl{part1-vocabulary} 资产为\textbf{已迁移的术语子集}建立可
查询锚点，保存词条号、标题、证据坐标和受控来源标识。以下 item 片段是该数据形态的
\Needspace{8\baselineskip}
书面投影；不得从本书代码块复制出第二份运行资产：

\nopagebreak[4]
\begin{codebox}
\begin{Verbatim}[fontsize=\small,breaklines=true,breakanywhere=true,breaksymbolleft={},breaksymbolright={}]
iso262:Term_1_item a iso262:Clause ;
    iso262:clauseId "1-3.84" ;
    dcterms:title "item"@en ;
    iso262:sourcePdfPage 24 ; iso262:sourceBlock 5 ;
    iso262:sourceBbox "58,193,102,206" .
iso262:Item iso262:hasSourceAnchor iso262:Term_1_item .
\end{Verbatim}
\end{codebox}

第二条三元组把 \nolinkurl{Item} 类挂到词条锚点上。Semantica 场景可以检查这条绑定及已编码
坐标合同；至于 \nolinkurl{Item} 的建模是否完整忠实于 1-3.84，仍须回到合法持有的原文和建模
语境人工审查。坐标可解析也不表示 OCR、中文转述或原 PDF 语义已经被机器核对。

下面的 SPARQL 只解释 package 查询想回答的问题；实际查询文本与 oracle 的唯一正本在
\Needspace{8\baselineskip}
Semantica normative package：

\nopagebreak[4]
\begin{codebox}
\begin{Verbatim}[fontsize=\small,breaklines=true,breakanywhere=true,breaksymbolleft={},breaksymbolright={}]
SELECT ?term ?clauseId ?title WHERE {
  ?cls iso262:hasSourceAnchor ?term .
  ?term iso262:clauseId ?clauseId ; dcterms:title ?title .
} ORDER BY ?clauseId
\end{Verbatim}
\end{codebox}

这条查询不会返回本附录的五个概念群，也不会给出首讲章，更不能再生 C.7.2 的字母索引：当前 normative package 的已迁移 RDF 资产没有保存这些编辑属性。C.1–C.5 的分群、C.7.1 的首讲速查、C.7.2 的字母索引以及各章正文回指，仍是人工维护的阅读视图。锚点给它们提供统一的受控标识和抽查入口，可以降低漂移风险，但尚未做到“改一处即全网一致”。修订时必须同时审查 Semantica 绑定与这些 Markdown 视图。

当前自动检查的能力边界如下。\nolinkurl{semantica.vol2.normative.scenario.modality-fidelity}
只执行其 CQ registry 明示的机器合同；坐标或绑定通过不表示 OCR、词条转述或原 PDF
已经完成语义核对，它也不读取本附录 Markdown，因而发现不了正文写错词条号。
\nolinkurl{severity} 等 OCR 噪声仍须回看原 PDF 并以校订留痕处理，不能归功于绑定查询。

诚实边界也要说清：Part 1 的出版覆盖分母始终是 185，实际机器锚定数只能从源锁定
Semantica manifest/receipt 查询，不能硬编码成永久进度。未锚定词条可以在回看原文后
用于正文，但不能假装已有机器门禁保护。若未来让分群、索引和首讲回指由查询生成，
必须先把编辑属性注册进 Semantica package，并为生成结果与 Markdown 建立独立校验；
在那之前仍需人工编辑审查。


\Needspace{6\baselineskip}
\section{字母索引、首讲章与交叉引用}

\Needspace{5\baselineskip}
\subsection{首讲章速查}

\begin{center}
\begingroup\small
\setlength{\tabcolsep}{4pt}
\begin{tabularx}{\linewidth}{@{}XXX@{}}
\toprule
概念群 & 首讲章 & 深讲/操作章 \\
\midrule
结构群（item/system/element/component/unit、SEooC） & ch02 & 各章；软件颗粒 ch07；半导体载体 appA \\
因果群（fault/error/failure 主链） & ch02 & \(\lambda\) 五分类 ch06；依赖失效与 FFI ch08 \\
风险群（HARA、S/E/C、ASIL、SG） & ch02 概览 & 操作 ch04；MSIL 换算 appB \\
需求群（SG\(\rightarrow\)FSR\(\rightarrow\)TSR、工作产物、确认措施） & ch02 概览 & ch04/ch05；档案与文化 ch03；工具与 PiU/DIA ch10 \\
时间群（FTTI 家族、安全状态） & ch02 概览 & 预算分配 ch05；现场时序 ch09 \\
\bottomrule
\end{tabularx}
\endgroup
\end{center}

\Needspace{5\baselineskip}
\subsection{字母序全量入口（185 条）}

以下按英文首字母列出全部 185 个受控词条的“词条号 · 本附录位置”。\textbf{粗体}为正文已展开条目（含合并讲解的成对条目），未加粗者见 C.7.3 速览。

\begin{itemize}
\item \textbf{A}：architecture \textbf{3.1}(C.1.2)；ASIL capability \textbf{3.2}(C.3.3)；ASIL decomposition \textbf{3.3}(C.3.3)；assessment \textbf{3.4}(C.4.4)；audit \textbf{3.5}(C.4.4)；automotive safety integrity level \textbf{3.6}(C.3.3)；availability 3.7(C.7.3)
\item \textbf{B}：base failure rate \textbf{3.8}(C.2.1)；base vehicle 3.9；baseline \textbf{3.10}(C.1.3)；body builder 3.11；body builder equipment 3.12；branch coverage \textbf{3.13}(C.4.5)；bus 3.14
\item \textbf{C}：calibration data \textbf{3.15}(C.1.3)；candidate \textbf{3.16}(C.1.3)；cascading failure \textbf{3.17}(C.2.4)；common cause failure \textbf{3.18}(C.2.4)；common mode failure \textbf{3.19}(C.2.4)；complete vehicle 3.20；component \textbf{3.21}(C.1.1)；configuration data \textbf{3.22}(C.1.3)；confirmation measure \textbf{3.23}(C.4.4)；confirmation review \textbf{3.24}(C.4.4)；controllability \textbf{3.25}(C.3.2)；coupling factors \textbf{3.26}(C.2.4)
\item \textbf{D}：dedicated measure \textbf{3.27}(C.4.2)；degradation \textbf{3.28}(C.5)；dependent failures \textbf{3.29}(C.2.4)；dependent failure initiator \textbf{3.30}(C.2.4)；detected fault \textbf{3.31}(C.2.3)；development interface agreement \textbf{3.32}(C.4.6)；diagnostic coverage \textbf{3.33}(C.2.5)；diagnostic points \textbf{3.34}(C.2.5)；diagnostic test time interval \textbf{3.35}(C.5)；distributed development \textbf{3.36}(C.4.6)；diversity \textbf{3.37}(C.4.2)；dual-point failure \textbf{3.38}(C.2.3)；dual-point fault \textbf{3.39}(C.2.3)
\item \textbf{E}：electrical and/or electronic system \textbf{3.40}(C.1.2)；element \textbf{3.41}(C.1.1)；embedded software \textbf{3.42}(C.1.1)；emergency operation \textbf{3.43}(C.5)；emergency operation time interval \textbf{3.44}(C.5)；emergency operation tolerance time interval \textbf{3.45}(C.5)；error \textbf{3.46}(C.2.1)；expert rider \textbf{3.47}(C.3.2)；exposure \textbf{3.48}(C.3.2)；external measure \textbf{3.49}(C.3.2)
\item \textbf{F}：failure \textbf{3.50}(C.2.1)；failure mode \textbf{3.51}(C.2.1)；failure mode coverage \textbf{3.52}(C.2.5)；failure rate \textbf{3.53}(C.2.1)；fault \textbf{3.54}(C.2.1)；fault detection time interval \textbf{3.55}(C.5)；fault handling time interval \textbf{3.56}(C.5)；fault injection \textbf{3.57}(C.4.2)；fault model \textbf{3.58}(C.2.1)；fault reaction time interval \textbf{3.59}(C.5)；fault tolerance \textbf{3.60}(C.2.3)；fault tolerant time interval \textbf{3.61}(C.5)；field data \textbf{3.62}(C.4.6)；formal notation \textbf{3.63}(C.4.5)；formal verification \textbf{3.64}(C.4.5)；freedom from interference \textbf{3.65}(C.2.4)；functional concept \textbf{3.66}(C.1.3)；functional safety \textbf{3.67}(C.3.1)；functional safety concept \textbf{3.68}(C.4.1)；functional safety requirement \textbf{3.69}(C.4.1)
\item \textbf{H}：hardware architectural metrics \textbf{3.70}(C.2.5)；hardware part \textbf{3.71}(C.1.1)；hardware elementary subpart \textbf{3.72}(C.1.1)；hardware subpart \textbf{3.73}(C.1.1)；harm \textbf{3.74}(C.3.1)；hazard \textbf{3.75}(C.3.1)；hazard analysis and risk assessment \textbf{3.76}(C.3.1)；hazardous event \textbf{3.77}(C.3.1)
\item \textbf{I}：independence \textbf{3.78}(C.2.4)；independent failures \textbf{3.79}(C.2.4)；informal notation \textbf{3.80}(C.4.5)；inheritance \textbf{3.81}(C.4.1)；inspection \textbf{3.82}(C.4.4)；intended functionality \textbf{3.83}(C.1.3)；item \textbf{3.84}(C.1.1)
\item \textbf{L}：latent fault \textbf{3.85}(C.2.3)；lifecycle \textbf{3.86}(C.4.3)
\item \textbf{M}：management system \textbf{3.87}(C.4.3)；malfunctioning behaviour \textbf{3.88}(C.2.1)；maximum time to repair time interval \textbf{3.89}(C.5)；model-based development \textbf{3.90}(C.4.2)；modification 3.91；modified condition/decision coverage \textbf{3.92}(C.4.5)；motorcycle 3.93；motorcycle safety integrity level \textbf{3.94}(C.3.3)；multi-core \textbf{3.95}(C.1.2)；multiple-point failure \textbf{3.96}(C.2.3)；multiple-point fault \textbf{3.97}(C.2.3)；multiple-point fault detection time interval \textbf{3.98}(C.5)
\item \textbf{N}：new development 3.99；non-functional hazard \textbf{3.100}(C.3.1)
\item \textbf{O}：observation points \textbf{3.101}(C.2.5)；operating mode \textbf{3.102}(C.3.1)；operating time \textbf{3.103}(C.5)；operational situation \textbf{3.104}(C.3.1)；other technology \textbf{3.105}(C.1.2)
\item \textbf{P}：partitioning \textbf{3.106}(C.1.2)；passenger car 3.107；perceived fault \textbf{3.108}(C.2.3)；permanent fault \textbf{3.109}(C.2.2)；phase \textbf{3.110}(C.4.3)；physics of failure 3.111；power take-off 3.112；processing element \textbf{3.113}(C.1.2)；programmable logic device \textbf{3.114}(C.1.2)；proven in use argument \textbf{3.115}(C.4.6)；proven in use credit \textbf{3.116}(C.4.6)
\item \textbf{Q}：quality management \textbf{3.117}(C.3.3)
\item \textbf{R}：random hardware failure \textbf{3.118}(C.2.2)；random hardware fault \textbf{3.119}(C.2.2)；reasonably foreseeable \textbf{3.120}(C.3.2)；rebuilding 3.121；redundancy \textbf{3.122}(C.4.2)；regression strategy \textbf{3.123}(C.4.4)；remanufacturing 3.124；residual fault \textbf{3.125}(C.2.3)；residual risk \textbf{3.126}(C.3.2)；review \textbf{3.127}(C.4.4)；risk \textbf{3.128}(C.3.2)；robust design \textbf{3.129}(C.4.2)
\item \textbf{S}：safe fault \textbf{3.130}(C.2.3)；safe state \textbf{3.131}(C.5)；safety \textbf{3.132}(C.3.1)；safety activity \textbf{3.133}(C.4.3)；safety anomaly \textbf{3.134}(C.2.4)；safety architecture \textbf{3.135}(C.1.2)；safety case \textbf{3.136}(C.4.3)；safety culture \textbf{3.137}(C.4.3)；safety element out of context \textbf{3.138}(C.1.3)；safety goal \textbf{3.139}(C.4.1)；safety manager \textbf{3.140}(C.4.3)；safety measure \textbf{3.141}(C.4.2)；safety mechanism \textbf{3.142}(C.4.2)；safety plan \textbf{3.143}(C.4.3)；safety-related element \textbf{3.144}(C.1.3)；safety-related function \textbf{3.145}(C.1.3)；safety-related incident 3.146；safety-related special characteristic \textbf{3.147}(C.4.6)；safety validation \textbf{3.148}(C.4.4)；semi-formal notation \textbf{3.149}(C.4.5)；semi-formal verification \textbf{3.150}(C.4.5)；semi-trailer 3.151；series production road vehicle 3.152；service note 3.153；severity \textbf{3.154}(C.3.2)；single-point failure \textbf{3.155}(C.2.3)；single-point fault \textbf{3.156}(C.2.3)；software component \textbf{3.157}(C.1.1)；software tool \textbf{3.158}(C.4.6)；software unit \textbf{3.159}(C.1.1)；statement coverage \textbf{3.160}(C.4.5)；sub-phase \textbf{3.161}(C.4.3)；supply agreement \textbf{3.162}(C.4.6)；system \textbf{3.163}(C.1.1)；systematic failure \textbf{3.164}(C.2.2)；systematic fault \textbf{3.165}(C.2.2)
\item \textbf{T}：target environment \textbf{3.166}(C.1.2)；technical safety concept \textbf{3.167}(C.4.1)；technical safety requirement \textbf{3.168}(C.4.1)；testing \textbf{3.169}(C.4.4)；tractor 3.170；trailer 3.171；transducer \textbf{3.172}(C.1.2)；transient fault \textbf{3.173}(C.2.2)；truck 3.174；T\&B vehicle configuration 3.175
\item \textbf{U}：unreasonable risk \textbf{3.176}(C.3.2)
\item \textbf{V}：variance in T\&B vehicle operation 3.177；vehicle function \textbf{3.178}(C.1.3)；vehicle operating state \textbf{3.179}(C.3.1)；verification \textbf{3.180}(C.4.4)；verification review \textbf{3.181}(C.4.4)
\item \textbf{W}：walk-through \textbf{3.182}(C.4.4)；warning and degradation strategy \textbf{3.183}(C.5)；well-trusted \textbf{3.184}(C.4.6)；work product \textbf{3.185}(C.4.3)
\end{itemize}

\Needspace{5\baselineskip}
\subsection{未展开词条速览（车辆类别与 T\&B 领域等）}

以下 23 条未在正文展开，多为车辆类别与卡客车（T\&B）领域词汇，工程语境见附录 B；此处各给一句定位，词条号仍以 Part 1 §3 为准。

\begin{itemize}
\item \textbf{可用性（availability, 3.7）}——产品在给定条件下按需提供规定功能的能力；与安全相关但不是安全本身。
\item \textbf{基础车辆（base vehicle, 3.9）}——加装上装设备之前的 OEM 卡客车配置。
\item \textbf{上装企业（body builder, 3.11）}——在基础车辆上加装车身、货箱或设备的组织。
\item \textbf{上装设备（body builder equipment, 3.12）}——装到 T\&B 基础车辆上的机具、车身或货箱。
\item \textbf{客车（bus, 3.14）}——设计用于载人及行李、座位数九座以上的机动车。注意：此处是“公共汽车”，\textbf{不是数据总线}——这是本词表最容易望文生义的一条。
\item \textbf{完整车辆（complete vehicle, 3.20）}——装好上装设备的 T\&B 基础车辆整车，如垃圾收集车、自卸车。
\item \textbf{改型（modification, 3.91）}——从既有相关项创建新相关项；变更管理与在用证明都会引用它（深讲 ch10）。
\item \textbf{摩托车（motorcycle, 3.93）}——两轮或整备质量不超 800 kg 的三轮机动车（不含轻便摩托车），Part 12 的主角（见 appB）。
\item \textbf{全新开发（new development, 3.99）}——创建具有先前未规定功能的相关项/元素，或对既有功能的全新实现（对照 modification，深讲 ch10）。
\item \textbf{乘用车（passenger car, 3.107）}——以载人为主、含驾驶员不超九座的车辆。
\item \textbf{失效物理（physics of failure, 3.111）}——基于失效机理研究的可靠性科学方法（半导体语境见 appA）。
\item \textbf{取力器（power take-off, 3.112）}——让卡车/牵引车动力源驱动外接设备的接口。
\item \textbf{改装（rebuilding, 3.121）}——改变 T\&B 原始配置以执行不同任务。
\item \textbf{再制造（remanufacturing, 3.124）}——T\&B 服役一段时间后按原规格拆解并换新/翻新部件。
\item \textbf{安全相关事件（safety-related incident, 3.146）}——安全相关失效的实际发生；现场监控的记账对象（见 ch09）。
\item \textbf{半挂车（semi-trailer, 3.151）}／\textbf{挂车（trailer, 3.171）}／\textbf{牵引车（tractor, 3.170）}／\textbf{卡车（truck, 3.174）}——T\&B 车辆类别四件套：卡车运货，牵引车拖半挂，半挂以鞍座连接并把垂直载荷压在牵引车上，挂车则自承全部重量。
\item \textbf{量产道路车辆（series production road vehicle, 3.152）}——预期用于公共道路、非原型的车辆。
\item \textbf{维修提示（service note, 3.153）}——执行维修程序时须考虑的安全信息文档。
\item \textbf{T\&B 车辆配置（T\&B vehicle configuration, 3.175）}——运行中不变的基础车辆与上装设备技术特征，如轴距、轴荷分配。
\item \textbf{T\&B 运行差异（variance in T\&B vehicle operation, 3.177）}——服役期内因载货或牵引而动态特性不同的使用形态。
\end{itemize}

\Needspace{5\baselineskip}
\subsection{最易混的八对近邻}

按五群走一遍之后，把全表最容易在评审桌上引起争论的近邻对收拢在一处，作为自测清单：

\begin{enumerate}
\item \textbf{fault / failure}——异常状态 vs 行为终止；故障可以潜伏，失效是显现的结果（C.2.1）。
\item \textbf{hazard / hazardous event}——潜在伤害来源 vs 它与运行场景的组合；HARA 按后者逐行打分（C.3.1）。
\item \textbf{safety measure / safety mechanism}——活动或技术方案的全集 vs 做进产品里的技术方案；评审是措施但不是机制（C.4.2）。
\item \textbf{verification / validation}——对着规格问“做得对不对” vs 对着安全目标问“目标对不对、真达到没有”（C.4.4）。
\item \textbf{confirmation review / audit / assessment}——分别检查工作产物、过程、达成的功能安全，对象不同不可互相顶替（C.4.4）。
\item \textbf{FTTI / FHTI}——相关项的容忍时间 vs 机制的处理用时；二者属于不同对象，Part 1 词条注记用 FDTI + FRTI < 相关 FTTI 解释及时覆盖，项目义务仍须追到适用条款与安全要求（C.5）。
\item \textbf{single-point fault / residual fault}——都直接违背安全目标，分界只在该元素有没有安全机制覆盖（C.2.3）。
\item \textbf{ASIL / MSIL / QM}——汽车等级、摩托车等级（须换算）、以及根本“不是一个 ASIL”的质量管理档（C.3.3）。
\end{enumerate}

\Needspace{5\baselineskip}
\subsection{使用提醒}

查词的读者从 C.7.2 进，按字母找到词条号与所在小节；学概念的读者从 C.1–C.5 进，沿着五条路径把邻居一起带走；做符合性工作的读者请记住三级效力：本附录白话 < 首讲章的三步走精读 < 原文词条。\textbf{词条号与标准版次}是三级之间的主要回查键之一；跨版次、勘误或上下文变化时仍须核对定义正文与适用范围。最后重复一遍 C.0 的纪律：白话引路，原文裁决——引路人不上裁判席。
